\documentclass[11pt]{article}

\usepackage[T1]{fontenc}
\usepackage[utf8]{inputenc}
\usepackage{mathpazo}
\usepackage{amsmath,amssymb,amsthm}
\usepackage{textcomp}

\usepackage{geometry}
\usepackage{microtype}
\usepackage{setspace}
\usepackage{titlesec}
\usepackage{hyperref}
\usepackage{booktabs}
\usepackage{array}
\usepackage{tabularx}
\usepackage{enumitem}

\geometry{
  paper=a4paper,
  top=1.25in,
  bottom=1.25in,
  left=1.4in,
  right=1.4in
}

\onehalfspacing

\hypersetup{
  colorlinks=true,
  linkcolor=black,
  citecolor=black,
  urlcolor=black
}

\titleformat{\section}{\large\bfseries}{\thesection.}{0.6em}{}
\titleformat{\subsection}{\normalsize\bfseries}{\thesubsection.}{0.5em}{}

% Theorem environments
\newtheoremstyle{rthm}
  {10pt}{10pt}{\itshape}{0pt}{\bfseries}{.}{0.5em}{}
\theoremstyle{rthm}
\newtheorem{theorem}{Theorem}[section]
\newtheorem{corollary}[theorem]{Corollary}
\newtheorem{proposition}[theorem]{Proposition}

\newtheoremstyle{rdef}
  {8pt}{8pt}{\normalfont}{0pt}{\bfseries}{.}{0.5em}{}
\theoremstyle{rdef}
\newtheorem{definition}[theorem]{Definition}
\newtheorem{example}[theorem]{Example}

\theoremstyle{remark}
\newtheorem*{remark}{Remark}
\newtheorem*{interpretation}{Interpretation}

% Macros
\newcommand{\M}{\mathcal{M}}
\newcommand{\A}{\mathcal{A}}
\newcommand{\Om}{\Omega}
\newcommand{\RD}[1][\theta]{R_{D}^{#1}}
\newcommand{\Lam}[1][\theta]{\Lambda^{#1}}
\newcommand{\Pers}{\operatorname{Pers}}
\newcommand{\Dep}{\operatorname{Depth}}

\title{\textbf{The Persistence Hierarchy}\\[6pt]
\large Adjective Order as Reachability Gradient}
\author{Flyxion\\[4pt]
\small Independent Researcher}
\date{June 2026}

\begin{document}

\maketitle
\thispagestyle{empty}

\begin{abstract}
\noindent
Every native English speaker knows that \emph{a small red wooden box} sounds
right while \emph{a wooden red small box} sounds wrong. Almost none can
explain why. The standard linguistic account lists a fixed sequence---opinion,
size, age, shape, colour, origin, material, purpose, noun---but offers no
principled explanation of why this sequence and not another. This paper
proposes that adjective order encodes a \emph{persistence hierarchy}: the
ordering reflects how deeply each property is coupled to the identity and
continued existence of the object it modifies. Properties that survive more
perturbations of the object sit closer to the noun; properties that are
observer-dependent, contextual, or easily altered sit further away. This
account derives the standard ordering from a single explanatory principle
rather than memorising a chart. It connects naturally to the Spanish
ser/estar distinction, which grammaticalises the same hierarchy through
verb choice rather than word order, and to a broader architectural
principle visible in writing systems: that stable symbolic substrates
support low-cost refinements, with high-value distinctions encoded closest
to the base and low-value annotations pushed outward. The Arabic script's
consonantal skeleton and diacritic layer instantiate this logic in
orthographic form; the English adjective sequence instantiates it in
positional order. A cross-linguistic analysis of Arabic nominal modification
further reveals that the hierarchy connects to the classical Arabic
philosophical distinction between \textit{jawhar} (substance, essence)
and \textit{carad} (accident, contingent property)---a distinction developed
by Al-Farabi and Ibn Sina in which the English adjective-order rule
participates implicitly. The persistence hierarchy is further interpretable
as a reachability gradient in the sense of admissibility theory: properties
closer to the noun have higher reachability divergence, meaning that losing
the distinction they mark forecloses more futures for agents who depend on
accurate object representation. Native speakers can navigate the persistence
hierarchy without knowing the chart exists because they are tracking a
genuine feature of the world---the differential stability of object
properties---rather than applying an arbitrary syntactic convention.
\end{abstract}

\newpage
\setcounter{page}{1}

%=====================================================================
\section{The Chart Nobody Knows}
\label{sec:chart}
%=====================================================================

Ask a native English speaker whether \emph{a small red wooden box} sounds
natural. They will say yes immediately. Ask whether \emph{a wooden red small
box} sounds natural. They will say no, and they will feel the wrongness
viscerally, the way a wrong note in a familiar melody feels wrong.

Now ask them why.

Almost none can answer. They know the rule in the sense that they reliably
follow it, but they do not know that they know it. The rule has never been
explicitly taught to most native speakers. It is acquired implicitly, the way
children acquire the phonology of their language: through exposure, without
instruction, and without conscious access to the underlying principle.

The rule in question is the English adjective order constraint. When multiple
adjectives modify a noun, they must appear in a fixed sequence:

\medskip
\begin{center}
\renewcommand{\arraystretch}{1.3}
\begin{tabular}{r l l}
\toprule
\textbf{Position} & \textbf{Category} & \textbf{Examples} \\
\midrule
1 & Determiner   & a, the, this, those \\
2 & Quantity     & four, some, many \\
3 & Opinion      & adorable, awful, delicious \\
4 & Size         & big, tiny, massive \\
5 & Age          & old, young, ancient \\
6 & Shape        & round, square, oval \\
7 & Colour       & red, blue, faded \\
8 & Origin       & Italian, eastern, Greek \\
9 & Material     & wooden, silk, metal \\
10 & Purpose     & dining, sleeping, racing \\
  & \textbf{Noun} & box, table, car \\
\bottomrule
\end{tabular}
\end{center}
\medskip

This ordering is not a stylistic preference. It is a grammatical constraint.
Violations are not merely inelegant; they are ungrammatical in the same sense
that subject-verb disagreement is ungrammatical. A native speaker who
produces \emph{a wooden red small box} has made an error, even if they cannot
identify it as such when asked.

The standard linguistic account describes the sequence but does not explain
it. Why should opinion precede size? Why should material precede purpose but
follow origin? Why should colour sit between shape and origin rather than
elsewhere? The chart is a description of a pattern. It is not a theory of why
the pattern exists.

This paper proposes such a theory.

%=====================================================================
\section{The Persistence Hierarchy}
\label{sec:persistence}
%=====================================================================

\subsection{The basic idea}

The ordering of adjectives in English reflects the degree to which each
property is coupled to the identity and continued existence of the object it
modifies. Properties that are constitutive of the object---that survive many
perturbations without the object ceasing to be itself---sit close to the noun.
Properties that are observer-relative, contextual, or easily altered without
changing the object's identity sit farther away.

Call this the \emph{persistence hierarchy}. Formally:

\begin{definition}[Persistence of a Property]
\label{def:persistence}
Let $O$ be an object and $P$ a property attributed to $O$. The
\emph{persistence} of $P$ with respect to $O$ is the measure of the set of
perturbations of $O$ under which $P$ remains true:
\[
\Pers(P, O) = \mu\!\left(\{\phi \in \Phi : P(\phi(O)) = \mathbf{true}\}\right),
\]
where $\Phi$ is the space of possible perturbations of $O$ and $\mu$ is a
measure on $\Phi$ reflecting the prior probability of each perturbation.
\end{definition}

\begin{definition}[Persistence Depth]
\label{def:depth}
The \emph{persistence depth} of adjective category $\mathcal{C}$ is the
average persistence of properties in $\mathcal{C}$ across objects and domains:
\[
\Dep(\mathcal{C}) = \mathbb{E}_{O, P \in \mathcal{C}}\!\left[\Pers(P, O)\right].
\]
\end{definition}

The Persistence Hierarchy Hypothesis is then:

\begin{proposition}[Persistence Hierarchy]
\label{prop:hierarchy}
Adjectives in English are ordered such that adjective category $\mathcal{C}_1$
precedes $\mathcal{C}_2$ in the prenominal sequence if and only if
$\Dep(\mathcal{C}_1) < \Dep(\mathcal{C}_2)$.
\end{proposition}

In plain language: the less persistent the property, the farther from the noun.
The more persistent the property, the closer to the noun.

\subsection{Deriving the standard ordering}

Let us work through the standard sequence and verify that persistence depth
increases monotonically from opinion to noun.

\paragraph{Opinion (position 3).}
Opinion adjectives encode the speaker's evaluative response to the object.
\emph{Beautiful}, \emph{adorable}, \emph{awful}. These are maximally
observer-relative: the same object is beautiful to one observer and ugly to
another. The same observer may revise their opinion tomorrow. Opinion
properties are highly sensitive to perturbations of the \emph{observer}
rather than the object. Their persistence with respect to the object itself is
minimal. Hence they sit furthest from the noun.

\paragraph{Size (position 4).}
Size is more objective than opinion but highly context-relative.
\emph{Big} relative to what? A big mouse is smaller than a small elephant.
Size also changes: objects shrink, expand, are cut down or built up. The
property is less observer-relative than opinion but more mutable than
material constitution. Persistence is low-to-moderate.

\paragraph{Age (position 5).}
Age is more stable than size but not fully constitutive. An old table was
once a young table. The same object passes through age predicates over time.
Age is a temporal property that the object accumulates rather than a
structural property that defines it. Persistence is moderate: the property
will change, but more slowly than opinion or contextual size.

\paragraph{Shape (position 6).}
Shape is more stable than age in many cases. The round table remains round
across many perturbations---moving it, painting it, ageing it. Shape changes
require structural intervention. Persistence increases.

\paragraph{Colour (position 7).}
Colour sits between shape and origin. It is more stable than opinion (there
is an objective fact about the colour of the table) but less stable than
material (you can repaint the table without changing what it is made of).
Repainting is a less drastic intervention than reshaping, and reshaping is
less drastic than replacing the material entirely. The persistence ordering
follows: shape $<$ colour $<$ material is wrong; the correct ordering is
shape slightly more stable, then colour, then origin and material. The
chart confirms this: shape (6) precedes colour (7) which precedes origin (8)
and material (9).

\paragraph{Origin (position 8).}
Origin is historically fixed. An Italian table was made in Italy and cannot
retrospectively have been made in France. Origin is a relational property
fixed at creation. It cannot be changed by any perturbation of the object
that leaves the object in existence. Persistence is high.

\paragraph{Material (position 9).}
Material is constitutive in the strongest sense. A wooden table is made of
wood throughout its existence as that table. You can paint it, scratch it,
move it across the world---it remains wooden. To change the material is to
destroy the table and create a different object. Material properties survive
almost all perturbations short of destruction. Persistence is near-maximal.

\paragraph{Purpose (position 10).}
Purpose sits adjacent to the noun because it specifies what the object is
\emph{for}, which is intimately connected to what kind of thing it is. A
dining table is not merely a table that happens to be used for dining: the
purpose is built into the categorisation. Purpose properties are relational
but stable: a dining table does not cease to be a dining table when it is
not in use. Persistence is high, reflecting the tight coupling between
function and identity.

\paragraph{The noun (position 11).}
The noun is the maximally persistent residue: it is the compressed identity
of the object itself. All other properties are modulations of this core.

\subsection{Why the ordering is felt rather than known}

Native speakers acquire the persistence hierarchy implicitly because it tracks
a genuine feature of the world rather than an arbitrary convention. The
hierarchy of property-persistence is available to any cognitive system that
represents objects over time and across perturbations. Children who have
learned that wooden things remain wooden when painted, but beautiful things
can stop being beautiful without any physical change, have implicitly acquired
the persistence ordering of those two property types. The grammatical rule
is downstream of this cognitive representation, not upstream of it.

This explains an otherwise puzzling asymmetry: children acquire adjective
order before they can explain it, ESL learners can explain it before they
acquire it fluently. The native speaker is navigating the persistence
hierarchy; the ESL learner is consulting a map of it. Both arrive at the
same sequence, but by different routes, and the native speaker's route is
faster and more reliable in production because it is grounded in object
cognition rather than metalinguistic rule-following.

%=====================================================================
\section{The Spanish Parallel: Ser and Estar}
\label{sec:spanish}
%=====================================================================

Spanish provides independent evidence for the persistence hierarchy by
grammaticalising it through a different means: the choice between two copulas.

English uses a single verb \emph{to be} for all predications. Spanish uses
two: \emph{ser} and \emph{estar}. The traditional description contrasts
permanent properties (ser) with temporary ones (estar), but this description
is incomplete and generates famous counterexamples. A better description is
that ser encodes properties treated as constitutive or identity-defining, while
estar encodes properties treated as contingent, resultant, or context-specific.

\medskip
\begin{center}
\renewcommand{\arraystretch}{1.3}
\begin{tabular}{p{5cm} p{5cm}}
\toprule
\textbf{Ser} (constitutive) & \textbf{Estar} (contingent) \\
\midrule
La mesa \emph{es} de madera. & La mesa \emph{est\'a} rota. \\
The table is (made of) wood. & The table is broken. \\[4pt]
\'El \emph{es} m\'edico. & \'El \emph{est\'a} enfermo. \\
He is a doctor. & He is sick. \\[4pt]
La silla \emph{es} redonda. & La silla \emph{est\'a} sucia. \\
The chair is round. & The chair is dirty. \\[4pt]
El cielo \emph{es} azul. & El cielo \emph{est\'a} nublado. \\
The sky is blue (by nature). & The sky is cloudy (right now). \\
\bottomrule
\end{tabular}
\end{center}
\medskip

The ser/estar distinction is a grammaticalisation of persistence depth. Ser
selects for high-persistence properties: material constitution, professional
identity, geometric shape, natural colour. Estar selects for low-persistence
properties: broken states, illness, dirt, current atmospheric conditions.

The famous puzzle cases resolve naturally under this analysis. \emph{\'El
est\'a muerto} (he is dead, using estar) seems to violate the
permanent/temporary rule since death is permanent. But death is a resultant
state, the outcome of a process, and estar is the verb of resultant and
contextual states. The persistence-hierarchy account handles this without
stipulation.

Similarly, \emph{la fruta est\'a verde} (the fruit is green, not yet ripe)
uses estar even though colour is a relatively stable property. But here
the speaker is treating the greenness as a contingent stage in the fruit's
development---it will ripen---rather than as a constitutive colour property.
The choice of estar signals that the colour is being construed as low-
persistence in this context. The same property can be construed as high- or
low-persistence depending on speaker intent, and ser/estar grammaticalises
that construal. English adjective order lacks this flexibility: material
always precedes colour in English because the ordering is fixed in the
grammar rather than available for construal-sensitive adjustment.

This cross-linguistic comparison suggests that the persistence hierarchy is
a universal feature of human object cognition that different languages
grammaticalise in different ways. English encodes it as positional order.
Spanish encodes it as copula selection. Other languages may encode it
through other means: classifier systems, evidentiality hierarchies,
or nominal modification patterns.

%=====================================================================
\section{Adjective Order as Reachability Gradient}
\label{sec:reachability}
%=====================================================================

The persistence hierarchy connects naturally to the reachability framework
developed in a companion paper \cite{flyxion2026}. That paper introduces
\emph{reachability divergence} $\RD(X,Y)$ as the measure of the set of
futures accessible from meaning $X$ but not meaning $Y$ in domain $\theta$.
The central claim is that linguistic distinctions survive when collapsing
them would foreclose futures a community needs to keep open.

Adjective order can be interpreted as a spatial encoding of reachability
divergence. When an adjective $A$ with high persistence is placed close to
the noun, it occupies a position that signals: the distinction $A$/not-$A$
controls which object you are dealing with, and therefore controls which
futures are available to you. When an adjective $B$ with low persistence is
placed far from the noun, it signals: the distinction $B$/not-$B$ is
peripheral to the object's identity and controls relatively few futures.

More precisely:

\begin{definition}[Adjectival Reachability Divergence]
\label{def:ard}
Let $O$ be an object, $P$ a property, and $\theta$ a domain of action.
The \emph{adjectival reachability divergence} of $P$ with respect to $O$
in $\theta$ is
\[
\RD[\theta](P, O)
= \mu^\theta\!\left(\A^\theta_{O \models P} \;\triangle\; \A^\theta_{O \not\models P}\right),
\]
where $\A^\theta_{O \models P}$ is the admissible future region for an agent
who correctly knows that $O$ has property $P$, and $\A^\theta_{O \not\models P}$
is the admissible future region for an agent who correctly knows that $O$
lacks $P$.
\end{definition}

The Persistence Hierarchy Hypothesis then becomes a special case of a more
general reachability claim:

\begin{proposition}[Reachability Gradient]
\label{prop:gradient}
Adjective order in English encodes adjectival reachability divergence:
adjective category $\mathcal{C}_1$ precedes $\mathcal{C}_2$ in the prenominal
sequence if and only if $\RD[\theta](\mathcal{C}_1, O) < \RD[\theta](\mathcal{C}_2, O)$
on average across objects $O$ and domains $\theta$.
\end{proposition}

The connection to persistence depth is immediate: high-persistence properties
tend to have high reachability divergence because they are more constitutive
of the object's identity, and correct knowledge of constitutive properties
opens more futures than correct knowledge of peripheral properties.

Consider a concrete illustration. An agent navigating a woodworking domain
needs to know whether the material is wood or metal before selecting
tools, fasteners, finishing products, and joining methods. The distinction
wooden/non-wooden has high $\RD$ in this domain: knowing it correctly
opens the correct set of futures; not knowing it (or holding a merged
representation that does not mark it) forecloses those futures. By contrast,
knowing whether the object is \emph{beautiful} controls almost no futures
in the woodworking domain. The carpenter proceeds identically regardless.
The reachability divergence of \emph{beautiful}/not-\emph{beautiful} is
near zero in this domain.

The adjective-order rule places \emph{wooden} near the noun and
\emph{beautiful} far from it. The positional encoding tracks the
reachability significance of the distinction.

\subsection{Why violations feel wrong}

The reachability gradient account makes a specific prediction about
grammaticality judgements. Adjective-order violations should feel worse
when the misplaced adjective has higher reachability divergence. Specifically:

\begin{proposition}[Violation Severity]
\label{prop:severity}
The perceived ungrammaticality of an adjective-order violation is a
monotonically increasing function of the adjectival reachability divergence
of the displaced adjective.
\end{proposition}

This predicts that moving a material adjective inward (away from the noun
to a position that implies lower persistence) should produce stronger
ungrammaticality judgements than moving an opinion adjective outward (to a
position implying higher persistence than it merits). Informally:

\begin{center}
\begin{tabular}{lll}
\toprule
\textbf{Violation} & \textbf{Type} & \textbf{Predicted severity} \\
\midrule
a red wooden small box & colour before material misplaced & moderate \\
a wooden beautiful box & material before opinion & severe \\
a beautiful wooden box & opinion before material & mild (near-canonical) \\
a small wooden old box & size before material/age reordered & moderate \\
\bottomrule
\end{tabular}
\end{center}

The violation \emph{a wooden beautiful box} should be perceived as severely
ungrammatical because it places the high-reachability material adjective in
a low-reachability position, misrepresenting the persistence hierarchy of
the object's properties. The violation \emph{a beautiful wooden box}
is near the canonical order (opinion, then material) and should be perceived
as only mildly deviant or acceptable, because opinion naturally precedes
material in the standard sequence.

This prediction is testable against existing psycholinguistic data on
adjective order. Several studies have collected grammaticality ratings for
adjective-order violations \cite{whorf1956,martin1969,scott2002,cinque1994};
the reachability gradient account predicts a correlation between rated
ungrammaticality and the persistence depth of the displaced adjective.

\subsection{Implicit navigation without explicit knowledge}

The finding that native speakers feel adjective-order violations without
being able to explain them is predicted by the reachability account but
would be puzzling on a purely syntactic account. If adjective order were
an arbitrary grammatical convention, there would be no reason to feel it
rather than merely know it. Arbitrary conventions can be learned and applied
without affective response.

But if adjective order tracks the persistence hierarchy, and the persistence
hierarchy reflects genuine facts about how object properties interact with
future action, then violation of adjective order is a misrepresentation
of object structure. Agents who have learned to navigate the world using
objects as tools will have developed sensitivity to the persistence hierarchy
through ordinary cognition. The grammatical rule is felt because it encodes
something real that object cognition already tracks.

This is analogous to why mathematical proofs feel compelling to trained
mathematicians even before they have fully verified every step: the proof
tracks genuine structure, and the trained mind can detect that tracking
even partially. A grammaticality judgement on adjective order is a
compressed evaluation of whether the property ordering matches the
persistence hierarchy of the object being described.

%=====================================================================
\section{Nested Structure and the Spherepop Interpretation}
\label{sec:spherepop}
%=====================================================================

There is a natural interpretation of the persistence hierarchy in terms of
nested process structure. Consider the noun not as a static object but as
the compressed residue of a generative process: the table is the outcome
of selecting material, imposing shape, applying finish, assigning function,
and locating in a context.

Each adjective layer represents a shell in the process that generated the
object. Properties at inner shells are constitutive of the process outcome.
Properties at outer shells are features added after the constitutive process
is complete.

\medskip
\begin{center}
\begin{tabular}{cl}
\toprule
\textbf{Shell} & \textbf{Property type} \\
\midrule
Innermost & Material: constitutive of what the object is \\
& Origin: historically fixed at creation \\
& Purpose: built into the object's functional identity \\
\midrule
Middle & Colour: applied after shaping \\
& Shape: imposed on material \\
& Age: accumulated over time \\
\midrule
Outermost & Size: relational and context-dependent \\
& Opinion: observer-relative and mutable \\
\bottomrule
\end{tabular}
\end{center}
\medskip

In this picture, the noun is the innermost bubble in a nested structure.
Each adjective is a layer around that bubble. The layers are ordered by
how closely they participate in the generative process that produces the
object. Inner layers can be thought of as programming the object's
identity; outer layers annotate it from the outside.

The standard adjective-order sequence then reads as a traversal of this
nested structure from outside in: the reader or listener encounters the
most contingent, observer-relative properties first, then progressively
closer to the constitutive core, arriving at the noun as the innermost
identity.

This traversal order has a functional interpretation. In language processing,
early adjectives set up the context of evaluation and constrain the space
of possible referents coarsely. Later adjectives narrow the referent space
more finely. Opinion tells you roughly what evaluative class the object
falls into. Size narrows it further. By the time you reach material, you
are very close to uniquely identifying the object type. The noun then
closes the identification. The ordering is cognitively efficient: it
moves from high-entropy, context-setting information to low-entropy,
identity-fixing information.

This efficiency is not separable from the persistence hierarchy. It is
the same hierarchy viewed from two directions. Moving from observer to
object is simultaneously moving from contingent to constitutive, from
high-entropy to low-entropy, from mutable to persistent.

%=====================================================================
\section{Script Architecture and the Incremental Augmentation Principle}
\label{sec:script}
%=====================================================================

The persistence hierarchy in adjective order is not an isolated phenomenon.
It belongs to a broader family of structures in which a stable substrate
supports layered, low-cost refinements that increase informational resolution
without requiring wholesale redesign. Writing systems exhibit this architecture
with unusual clarity, and examining it strengthens the theoretical basis for
the persistence hierarchy account.

\subsection{Stable substrate and minimal refinement}

Consider how the Arabic consonantal script handles the problem of phonemic
distinction. Early manuscripts were written in a consonantal skeleton---the
\textit{rasm}---in which many letters that are today graphically distinct
shared identical base shapes. The forms now written as \textbf{b}, \textbf{t},
and \textbf{th} were graphically identical. Disambiguation relied on context.

Over time, a solution emerged that did not abandon the efficient skeletal
layer: systematic dotting. A single dot placed above or below the shared base
form was sufficient to differentiate multiple phonemes. The informational
return of this augmentation is disproportionately large relative to its cost.
A dot is a minimal motor gesture, yet it yields a categorical distinction.

The structural principle this illustrates---call it \emph{incremental
augmentation}---is that symbolic systems under constraint tend to introduce
new distinctions through minimal additions to existing structure rather than
through wholesale replacement. The base layer is preserved because it is
already efficient; the augmentation is accepted because its cost is low
relative to its informational benefit.

\begin{definition}[Incremental Augmentation]
\label{def:aug}
Let $X$ be a base symbolic structure and $A$ a minimal augmentation operator.
The augmentation is \emph{efficient} if

\[
C(A) \ll C(X)
\quad \text{while} \quad
I(A(X)) - I(X) \gg 0,
\]

where $C$ is production cost and $I$ is informational distinctness. An
efficient augmentation achieves high discriminative gain at low additional
effort.
\end{definition}

The dotted Arabic letters are a paradigm case. The base form $b$ (bare
skeletal shape) and the augmented form $d(b)$ (dot added above) satisfy
$C(d) \approx \varepsilon$ while $D(b, d(b)) \geq D_{\min}$: minimal cost,
maximal discriminative payoff.

\subsection{The adjective hierarchy as augmentation sequence}

The persistence hierarchy is a cognitive instance of the same principle.
The noun is the base form---the maximally compressed representation of the
object. Each adjective layer is an augmentation that adds information at
increasing cost as you move outward from the noun toward the observer.

Material augmentation (``wooden'') has high informational return because it
is constitutive: knowing it correctly opens the full set of futures associated
with the object type. Its informational leverage is high. Opinion augmentation
(``beautiful'') has low informational return because it is observer-relative:
knowing it correctly adds little to your navigational capacity relative to
knowing the material and origin. Its leverage is low.

The ordering opinion $\to$ material is therefore an ordering from
low-leverage augmentation to high-leverage augmentation as you approach the
noun. The noun is the base; the adjective sequence is a series of augmentations
arranged so that the highest-leverage refinements sit closest to what they
refine.

This is the same efficiency logic as the dot. Put the high-value distinctions
closest to the structure they modify. Put the low-value annotations further
away, where they can be dropped if context permits without losing navigational
capacity.

\subsection{Positional regularization as compression}

Writing systems that allow fully context-sensitive joining face a combinatorial
explosion: if the shape of each letter depends on both its left and right
neighbours, the number of possible forms grows as $\mathcal{O}(n^2)$ for
pairwise dependencies. Arabic solved this by reducing context-sensitive
variation to four predictable positional allographs---isolated, initial,
medial, and final---which scale as $\mathcal{O}(4n)$, linear rather than
quadratic. The combinatorial space is compressed by replacing local variation
with a global positional rule.

Adjective order performs exactly this compression for nominal modification.
Rather than permitting every possible ordering of adjectives and resolving
ambiguity pragmatically, English fixes the positions so that the parsing
load is eliminated. A native speaker does not need to consider whether
\emph{wooden red small} might mean something different from \emph{small red
wooden}: the positional encoding renders one of them ungrammatical and
removes the interpretive computation entirely. The grammar pays a small cost
in flexibility to achieve a large reduction in processing load---the same
trade as the script makes when it moves from context-sensitive ligatures to
positional allographs.

\subsection{Non-connecting letters as controlled interruption}

Arabic scripts face a second optimization problem: maximizing cursive
continuity reduces production cost (fewer pen lifts) but risks perceptual
collapse (letters become visually indistinguishable when merged into
uninterrupted flow). The solution is selective interruption: a small set of
non-connecting letters terminates the cursive sequence at exactly the points
where continued joining would erode discriminability. These letters function
as structured boundary markers, paying a local cost in motor continuity to
preserve global legibility.

Adjective order has an analogue. The constraint is not uniform: it applies
more rigidly to some category pairs than others, and native speakers show
graded judgements rather than categorical rejection across all violations.
The rigidity is highest precisely where the positional encoding matters most
for disambiguation---where without the fixed order, the persistence hierarchy
of the object would be misrepresented. Where it matters less (two
high-similarity-persistence categories adjacent to each other), flexibility
appears. The grammar interrupts only where interruption is necessary to
preserve the hierarchy, not uniformly throughout.

\subsection{The general pattern}

Across these parallels---augmentation efficiency, positional compression,
controlled interruption---the same structural logic recurs. Systems under
embodied constraint stabilize not by maximizing any single variable (pure
motor efficiency, pure discriminability, pure flexibility) but by achieving
local optima where small adjustments no longer improve the balance between
competing pressures. The persistence hierarchy in adjective order is one
such equilibrium: it is the ordering that maximizes navigational information
delivered per unit of positional slot, given the constraint that the slots
must be linearly ordered and that the noun anchors the right edge.

This suggests that the English adjective-order rule is not an accident of
grammatical history but a convergent solution to a constraint satisfaction
problem that writing systems, phonological systems, and nominal modification
systems all face in their own medium.

%=====================================================================
\section{Arabic and the Persistence Hierarchy}
\label{sec:arabic}
%=====================================================================

Arabic provides an independent and typologically distant illustration of the
persistence hierarchy. Egyptian Arabic, like Modern Standard Arabic, does not
grammaticalise adjective order as rigidly as English, but the hierarchy is
nonetheless visible in how speakers and analysts articulate the relationship
between adjectives and the nouns they modify.

Consider the following analysis, offered by a native Arabic speaker reflecting
on the English adjective-order phenomenon:

\medskip
\begin{quote}
\textit{el-fekra di acma' \v{s}wayya men e\v{s}-\v{s}arh et-taql\={i}di li-tart\={i}b es-sif\={a}t
fil-ingil\={i}zi.}\\
\smallskip
\textsc{the-idea this deeper a-little from the-explanation traditional
of-order the-adjectives in-English}\\
\smallskip
``This idea is a little deeper than the traditional explanation of adjective
order in English.''
\end{quote}

\begin{quote}
\textit{et-tafs\={i}r illi ana \v{s}\={a}yfo inn tart\={i}b es-sif\={a}t biyyackis daragit
irtib\={a}t es-sifa bihuwiyyit e\v{s}-\v{s}ay' nafsu.}\\
\smallskip
\textsc{the-explanation that I see-it that order the-adjectives reflects
degree attachment the-adjective to-identity the-thing itself}\\
\smallskip
``The explanation I see is that adjective order reflects the degree to which
an adjective is attached to the identity of the thing itself.''
\end{quote}

\begin{quote}
\textit{kull m\={a} k\={a}nit es-sifa a'rab li-tab\={i}cit e\v{s}-\v{s}ay' wi-gawharu, kull m\={a}
'irribit min el-ism.}\\
\smallskip
\textsc{every time was the-adjective closer to-nature the-thing
and-essence-its, every time it-came-closer to the-noun}\\
\smallskip
``The closer the adjective is to the nature and essence of the thing, the
closer it comes to the noun.''
\end{quote}

\begin{quote}
\textit{huwwa mi\v{s} biytabbi' q\={a}cida nahwiyya w\={a}cya, lakinnu biyitc\={a}mil mac
taba'\={a}t mukhtalifa min il-irtib\={a}t bi\v{s}-\v{s}ay' nafsu.}\\
\smallskip
\textsc{he not applies rule grammatical conscious, but-he deals with layers
different of attachment to-the-thing itself}\\
\smallskip
``He is not consciously applying a grammatical rule; rather, he is dealing
with different layers of attachment to the thing itself.''
\end{quote}
\medskip

Several things are worth noting in this formulation.

\paragraph{The gloss makes constituency visible.}
The interlinear Arabic translation makes explicit a feature of the persistence
hierarchy that English obscures: the morphological transparency of the
attachment relation. Arabic \textit{irtib\={a}t} (attachment, connection) combined
with \textit{huwiyyit e\v{s}-\v{s}ay'} (identity of the thing) gives the phrase
\textit{daragit irtib\={a}t es-sifa bihuwiyyit e\v{s}-\v{s}ay'}---literally, the degree
of the adjective's attachment to the identity of the thing. The construct state
(\textit{id\={a}fa}) in Arabic grammaticalises the possession-and-attachment
relation directly: \textit{huwiyyit e\v{s}-\v{s}ay'} is ``the thing's identity'' as a
syntactic unit. The persistence hierarchy is therefore expressible in Arabic
with greater morphological directness than in English.

\paragraph{The observer/object distinction is grammatically marked.}
The Arabic formulation distinguishes between \textit{ra'y il-mur\={a}'ib}
(the observer's opinion, literally ``view of the observer'') and
\textit{huwiyyit e\v{s}-\v{s}ay'} (the identity of the thing). This is not merely a
conceptual distinction; the Arabic morphosyntax makes it a syntactic one.
The possessive construction in each case foregrounds whose perspective is
operative. Opinion adjectives belong to the observer; constitutive adjectives
belong to the thing. The persistence hierarchy is already a distinction between
observer-relative and object-relative properties, and Arabic grammaticalises
this distinction at the morphological level.

\paragraph{Arabic adjectival agreement reflects persistence depth.}
In Arabic, adjectives agree with their nouns in gender, number, definiteness,
and case. The agreement morphology creates a tight binding between noun and
modifier. Critically, adjectives in Arabic follow the noun rather than
preceding it: \textit{el-tawla el-kha\v{s}abiyya} (the table the-wooden, ``the
wooden table''). This post-nominal position means that Arabic encodes the
hierarchy differently from English---not through prenominal ordering but
through the sequence of post-nominal modifiers when multiple adjectives appear.

Studies of Arabic nominal modification suggest that when multiple adjectives
follow a noun in Arabic, the ordering preferences mirror the persistence
hierarchy: constitutive modifiers (material, origin) tend to appear closer to
the noun while contingent modifiers (opinion, temporary state) appear farther
away. The hierarchy is the same; the directionality is reversed because the
noun precedes rather than follows its modifiers.

\paragraph{The verbal gloss reveals the process structure.}
The Arabic reflection uses the verb \textit{biyyackis} (reflects, mirrors) to
describe what adjective order does to the persistence hierarchy. The choice of
this verb is significant: a mirror is a faithful representation of structure,
not a stipulation. The implication is that adjective order is not imposing an
order on properties but revealing an order that already exists in the
relationship between properties and object identity. This is precisely the
claim of the persistence hierarchy account: the grammatical rule is downstream
of the cognitive structure, not constitutive of it.

\paragraph{The key Arabic term: \textit{jawhar}.}
The word \textit{jawhar} (essence, substance, jewel) appears in the phrase
\textit{tab\={i}cit e\v{s}-\v{s}ay' wi-gawharu} (the nature of the thing and its essence).
In classical Arabic philosophy, \textit{jawhar} is the technical term for
substance in the Aristotelian sense---the category of being that persists
through change and constitutes the identity of a thing. Its appearance here
is not accidental. The intuition that material and constitutive properties sit
closest to the noun because they are closest to the \textit{jawhar} of the
object is a direct inheritance from the Arabic philosophical tradition's
engagement with Aristotelian substance metaphysics.

This is a deeper connection than a cross-linguistic parallel. It suggests
that the persistence hierarchy in adjective order is the grammatical
encoding of a metaphysical distinction---between substance and accident,
\textit{jawhar} and \textit{carad}---that has been articulated in the Arabic
philosophical tradition since at least Al-Farabi and Ibn Sina. The English
adjective-order rule, which most English speakers have never analysed, encodes
a distinction that medieval Arabic philosophers developed elaborate theoretical
machinery to describe.

The ESL learner consulting a chart is memorising a rule whose philosophical
depth they may not suspect. The native speaker following it implicitly is
navigating a distinction between substance and accident that Aristotle,
translated into Arabic and refined across centuries of Islamic philosophy,
made one of the central problems of metaphysics.

%=====================================================================
\section{Cross-Linguistic Predictions}
\label{sec:crossling}
%=====================================================================

If the persistence hierarchy is a universal feature of human object cognition,
it should leave traces in languages other than English, even when those
languages do not grammaticalise adjective order as rigidly.

Several predictions follow:

\paragraph{Languages with flexible adjective order should show persistence
effects in neutral contexts.}
In languages where adjective order is grammatically unconstrained, speakers
should nevertheless prefer orders that match the persistence hierarchy in
unmarked, neutral utterances. Deviations from the hierarchy should be
interpretable as marked: an adjective moved outward from its canonical
position may signal that the speaker is treating the property as more
contingent than usual.

\paragraph{Copula systems should reflect the hierarchy.}
Languages with multiple copulas should align their copula distinctions with
the persistence hierarchy. Spanish ser/estar is one case. Igbo \emph{b\`u}/
\emph{d\`\i{}} makes a similar distinction. Arabic \emph{kana}/\emph{laysa}
have related distribution patterns. The prediction is that the high-persistence
side of any copula distinction should select for the property types that sit
closer to the noun in English adjective order.

\paragraph{Classifier systems should track material and shape.}
Languages with nominal classifiers typically classify by material and shape
rather than by opinion or size. Mandarin measure words, Japanese classifiers,
and Burmese classifiers all prioritise the constitutive properties that sit
innermost in the English adjective-order sequence. This is expected if
classifiers are grammaticalisations of the most identity-relevant property
dimensions.

\paragraph{Adjective-order rigidity should correlate inversely with other
persistence-encoding mechanisms.}
Languages that grammaticalise the persistence hierarchy through copula
selection, classifier systems, or nominal morphology should show less rigid
adjective order, because the hierarchy is already encoded elsewhere.
Languages that lack these mechanisms should show more rigid adjective
ordering as the primary grammatical encoding of the hierarchy.

These are testable predictions that connect the persistence hierarchy account
to the typological literature on adjective ordering, copula systems, and
nominal classification.

%=====================================================================
\section{Implications for Language Acquisition}
\label{sec:acquisition}
%=====================================================================

The persistence hierarchy account makes specific predictions about the
acquisition of adjective order that differ from a purely syntactic account.

On a syntactic account, adjective order is a language-specific parameter
that children must learn by exposure to their target language. The order
is arbitrary relative to the world and must be memorised. Errors should
be random with respect to property type: children should be as likely to
misorder material before opinion as opinion before material.

On the persistence hierarchy account, the ordering reflects a non-linguistic
cognitive structure that children develop through ordinary object interaction.
Children who have learned that wooden objects remain wooden when painted,
but beautiful objects stop being beautiful when damaged, have implicitly
acquired the persistence ordering of those property types. Grammatical
acquisition then consists in aligning this pre-existing cognitive hierarchy
with the language-specific positional encoding.

This predicts that children's errors in adjective order should be
\emph{non-random}: errors involving high-persistence adjectives displaced to
low-persistence positions should be rarer than the reverse, because the
cognitive hierarchy already privileges material and origin adjectives as
identity-constitutive. Children may more readily err by placing opinion
adjectives too close to the noun (treating a subjective evaluation as more
constitutive than it is) than by placing material adjectives too far from
the noun (treating a constitutive property as more contingent than it is).

The ESL learner presents the mirror case. They have the cognitive persistence
hierarchy from their native-language experience and they have the chart from
instruction. Their errors arise not from lacking the hierarchy but from
failing to automate the mapping between the hierarchy and positional order.
They must consciously consult the rule because the rule has not been
integrated into the object-cognition process that the native speaker navigates
implicitly. Fluent production requires not knowing the chart but being the
chart.

%=====================================================================
\section{Relationship to the Reachability Theory of Lexical Preservation}
\label{sec:connection}
%=====================================================================

This paper is a companion to \cite{flyxion2026}, which argues that linguistic
distinctions survive when collapsing them would foreclose futures a community
needs to keep open. The present paper extends this framework from the
survival of lexical distinctions to the internal ordering of nominal modifiers.

The connection runs in both directions.

\paragraph{Adjective order as implicit stability.}
The persistence hierarchy is a clear case of a linguistic structure that is
highly stable without being consciously maintained. No speaker enforces
adjective order by consulting the chart. The rule survives because it tracks
the persistence hierarchy, and the persistence hierarchy tracks genuine facts
about object constitution that agents need for navigation. This supports the
companion paper's claim that linguistic structures can be preserved by
navigational pressure rather than communicative convention or conscious
rule-following.

\paragraph{The persistence hierarchy as reachability ordering.}
Proposition~\ref{prop:gradient} establishes that adjective order is a
reachability gradient: position encodes reachability divergence. High-
persistence properties have high reachability divergence and sit close to
the noun. Low-persistence properties have low reachability divergence and
sit far from the noun. The positional encoding is therefore a spatial
representation of the same quantity---reachability divergence---that the
companion paper shows governs the survival of lexical distinctions.

\paragraph{A unified picture.}
Both phenomena---which lexical distinctions survive compression, and in what
order adjectives appear before a noun---are governed by the same underlying
quantity: how much future access depends on the distinction being preserved.
Lexical distinctions with high reachability divergence resist colexification.
Adjectival properties with high reachability divergence sit close to the noun.
Both are expressions of a single principle: language preserves and foregrounds
the distinctions that matter most for navigating the futures agents face.

This suggests that the reachability framework is not specific to lexical
organisation but extends to syntactic organisation as well. The persistence
hierarchy in adjective order is a syntactic implementation of the same
principle that governs lexical stability: put the high-stakes distinctions
where they cannot be lost.

%=====================================================================
\section{Conclusion}
\label{sec:conclusion}
%=====================================================================

The English adjective-order constraint is one of the most widely known
examples of implicit grammatical knowledge: a rule that native speakers
reliably follow without knowing it exists. The standard linguistic account
lists the sequence but does not explain it.

This paper proposes that the sequence encodes a persistence hierarchy: the
ordering from opinion to material reflects the ordering from least persistent
to most persistent among the properties an adjective can express. Properties
that survive more perturbations of the object, that are more constitutive of
its identity, sit closer to the noun. Properties that are observer-relative,
contextual, or easily altered sit further away.

The persistence hierarchy derives the standard adjective order from a single
explanatory principle rather than a memorised chart. It explains why native
speakers feel the order rather than merely knowing it: the order tracks a
genuine feature of object constitution that cognition already represents. It
explains the Spanish ser/estar distinction as a grammaticalisation of the
same hierarchy through copula choice rather than word order. The same
efficiency logic---stable base, high-value distinctions closest to the base,
low-value annotations pushed outward---governs the Arabic consonantal skeleton
and its diacritic layer, the compression of context-sensitive ligatures into
positional allographs, and the selective use of non-connecting letters to
interrupt cursive flow only where discriminability requires it. Adjective
order and writing system design are convergent solutions to the same
constraint satisfaction problem. The
analysis of Arabic nominal modification further reveals that the hierarchy
connects to the classical philosophical distinction between \textit{jawhar}
and \textit{carad}---between substance and accident---a distinction that
Al-Farabi and Ibn Sina developed into a rigorous theoretical framework
while reflecting on the same Aristotelian intuitions that underlie the
English adjective-order rule. The chart that English teachers draw on
whiteboards is a pedagogical rediscovery of a metaphysical principle that
medieval Islamic philosophy had already formalised. The account generates
testable predictions about violation severity, cross-linguistic copula
systems, classifier hierarchies, and acquisition errors.

The persistence hierarchy is further interpretable as a reachability gradient:
adjective position encodes how much future access depends on knowing the
property correctly. This connects adjective order to a companion theory of
lexical preservation, suggesting that both phenomena---which distinctions
survive compression and how modifiers are ordered---are governed by the same
underlying principle.

Language puts the things that matter most for navigation where they are
hardest to lose.

%=====================================================================
\begin{thebibliography}{99}

\bibitem{flyxion2026}
Flyxion (2026).
Why distinctions survive: A reachability theory of lexical preservation.
\textit{Independent Research Monograph}.

\bibitem{avicenna1027}
Ibn S\={\i}n\=a (Avicenna) (c.\ 1027).
\textit{Kit\=ab al-Shif\=a'} (The Book of Healing).
Partial English translation: A.\ M.\ Goichon (1938),
\textit{La distinction de l'essence et de l'existence d'apr\`es Ibn S\={\i}n\=a}.
Descl\'ee de Brouwer.

\bibitem{alfarabi930}
Al-F\=ar\=ab\={\i} (c.\ 930).
\textit{Kit\=ab al-Hur\=uf} (The Book of Letters).
Edition and commentary: M.\ Mahdi (1969).
Dar el-Mashreq.

\bibitem{cinque1994}
Cinque, G.\ (1994).
On the evidence for partial N-movement in the Romance DP.
In G.\ Cinque et al.\ (Eds.), \textit{Paths Towards Universal Grammar},
85--110. Georgetown University Press.

\bibitem{cinque2010}
Cinque, G.\ (2010).
\textit{The Syntax of Adjectives: A Comparative Study}.
MIT Press.

\bibitem{dixon1982}
Dixon, R.\ M.\ W.\ (1982).
\textit{Where Have All the Adjectives Gone?}
Mouton.

\bibitem{martin1969}
Martin, J.\ E.\ (1969).
Semantic determinants of preferred adjective order.
\textit{Journal of Verbal Learning and Verbal Behavior}, 8, 697--704.

\bibitem{murphy2002}
Murphy, G.\ L.\ (2002).
\textit{The Big Book of Concepts}.
MIT Press.

\bibitem{scott2002}
Scott, G.-J.\ (2002).
Stacked adjectival modification and the structure of nominal phrases.
In G.\ Cinque (Ed.), \textit{Functional Structure in DP and IP}, 91--120.
Oxford University Press.

\bibitem{svenonius2008}
Svenonius, P.\ (2008).
The position of adjectives and other phrasal modifiers in the decomposition
of DP.
In L.\ McNally \& C.\ Kennedy (Eds.),
\textit{Adjectives and Adverbs}, 16--42. Oxford University Press.

\bibitem{whorf1956}
Whorf, B.\ L.\ (1956).
\textit{Language, Thought, and Reality}.
MIT Press.

\end{thebibliography}

\end{document}
